\documentclass[10pt,a4paper]{article}
\usepackage[utf8]{inputenc}
\usepackage[T1]{fontenc}
\usepackage{amsmath}
\usepackage{amssymb}
\usepackage{graphicx}
\usepackage[spanish]{babel}
\usepackage{booktabs}
\usepackage{array}
\usepackage{placeins}

\usepackage{url}
\usepackage{hyperref}
\usepackage{listings}
\usepackage{textcomp}
\usepackage{accsupp}
\usepackage{attachfile}

% color def
\usepackage{color}
\definecolor{darkred}{rgb}{0.6,0.0,0.0}
\definecolor{darkgreen}{rgb}{0,0.50,0}
\definecolor{lightblue}{rgb}{0.0,0.42,0.91}
\definecolor{orange}{rgb}{0.99,0.48,0.13}
\definecolor{grass}{rgb}{0.18,0.80,0.18}
\definecolor{pink}{rgb}{0.97,0.15,0.45}

% listings
\usepackage{listings}

% General Setting of listings
\lstset{
	language=Python,
	tabsize=4,
	aboveskip=1em,
	breaklines=true,
	abovecaptionskip=-6pt,
	captionpos=b,
	escapeinside={\%*}{*)},
	frame=single,
	numbers=left,
	numbersep=8pt,
	numberstyle=\tiny,
	morekeywords=[1]{,as,assert,nonlocal,with,yield,self,True,False,None,},
	morekeywords=[2]{,__init__,__add__,__mul__,__div__,__sub__,__call__,__getitem__,__setitem__,__eq__,__ne__,__nonzero__,__rmul__,__radd__,__repr__,__str__,__get__,__truediv__,__pow__,__name__,__future__,__all__,},
	morekeywords=[3]{,object,type,isinstance,copy,deepcopy,zip,enumerate,reversed,list,set,len,dict,tuple,range,xrange,append,execfile,real,imag,reduce,str,repr,},
	morekeywords=[4]{,Exception,NameError,IndexError,SyntaxError,TypeError,ValueError,OverflowError,ZeroDivisionError,},
	morekeywords=[5]{,ode,fsolve,sqrt,exp,sin,cos,arctan,arctan2,arccos,pi, array,norm,solve,dot,arange,isscalar,max,sum,flatten,shape,reshape,find,any,all,abs,plot,linspace,legend,quad,polyval,polyfit,hstack,concatenate,vstack,column_stack,empty,zeros,ones,rand,vander,grid,pcolor,eig,eigs,eigvals,svd,qr,tan,det,logspace,roll,min,mean,cumsum,cumprod,diff,vectorize,lstsq,cla,eye,xlabel,ylabel,squeeze,},
	deletekeywords=[2]{next,set},
	deletekeywords=[3]{next,set},
	basicstyle=\ttfamily,
	backgroundcolor=\color{white},
	commentstyle=\color{darkgreen}\slshape,
	keywordstyle=\color{blue}\bfseries\itshape,
	keywordstyle=[2]\color{blue}\bfseries,
	keywordstyle=[3]\color{grass},
	keywordstyle=[4]\color{red},
	keywordstyle=[5]\color{orange},
	stringstyle=\color{darkred},
	emphstyle=\color{pink}\underbar,
	morestring=[d]{"},
	showstringspaces=false,
	upquote=true,
    literate={á}{{\'a}}1 {é}{{\'e}}1 {í}{{\'i}}1 {ó}{{\'o}}1 {ú}{{\'u}}1 {ñ}{{\~n}}1 {Á}{{\'A}}1 {É}{{\'E}}1 {Í}{{\'I}}1 {Ó}{{\'O}}1 {Ú}{{\'U}}1 {Ñ}{{\~N}}1 {¿}{{\textquestiondown}}1 {¡}{{\textexclamdown}}1
}

\usepackage{tikz}
\usetikzlibrary{calc,shapes.multipart,chains,arrows,positioning}

\makeatletter
\lst@RequireAspects{writefile}

% Use \attachfile command to add listing content as paperclip.
\lstnewenvironment{attachedlisting}[1][]{\lst@BeginAlsoWriteFile{#1.txt}% 
}
{% 
	\lst@EndWriteFile
	\marginpar{\attachfile[appearance=false,icon=Paperclip,mimetype=text/tex]{#1.txt}}
}

% Use accsupp package to add listing content as copyable text.
\lstnewenvironment{copyablelisting}{\lst@BeginAlsoWriteFile{\jobname.txt}% 
}
{% 
	\lst@EndWriteFile
	\let\verbatim@processline\add@lstline
	\global\let\lstfile\empty
	\verbatiminput{\jobname.txt}%
	\marginpar{(\BeginAccSupp{method=escape,ActualText={\lstfile}}\PaperPortrait\EndAccSupp{})}
}
\def\add@lstline
{\xdef\lstfile{\unexpanded\expandafter{\lstfile}\the\verbatim@line\string^^J}}
\makeatother

\tikzset{
	squarecross/.style={
		draw, rectangle,minimum size=18pt, fill=orange!80,
		inner sep=0pt, text=black,
		path picture = {
			\draw[black]
			(path picture bounding box.north west) --
			(path picture bounding box.south east)
			(path picture bounding box.south west) --
			(path picture bounding box.north east);
		}
	},
}

\title{Creando Nuestra Propia Lista Enlazada}
\date{}
\begin{document}
	\maketitle

\section*{¿Qué es una Lista Enlazada?}

Una \texttt{Lista Enlazada} (LinkedList) es una estructura de datos que funciona de manera similar a una \texttt{list} en Python, pero la forma en que los datos se almacenan en memoria es diferente. En una \texttt{LinkedList}, todos los elementos de la lista están enlazados o encadenados, permitiéndonos saltar de un elemento a otro.

En una \texttt{LinkedList} simplemente enlazada, cada elemento tendrá dos componentes. 
Uno contiene los datos que queremos almacenar en la lista y el otro contendrá una referencia al siguiente elemento de la lista. 
Usando esto podemos navegar de un elemento hacia adelante en la cadena, siempre en la misma dirección. 
Podemos detectar el último elemento ya que no referencia a ningún otro elemento. 
La forma en que los elementos están enlazados significa que si almacenamos la referencia al primer elemento en una variable, podremos llegar a cualquier otro elemento de la lista.
\vspace{0.7cm}
\begin{center}
	\shorthandoff{>}
	\begin{tikzpicture}[list/.style={
						very thick, rectangle split,
						rectangle split parts=2, draw,
						rectangle split horizontal, minimum size=18pt,
						inner sep=4pt, text=black,
						rectangle split part fill={red!20, blue!20}
					},
			->, start chain, very thick,
			chain_arrow/.style={*->, thick}]
		\node[list,on chain] (A) {cabeza};
		\node[list,on chain] (B) {medio};
		\node[list,on chain] (C) {cola};
		\node[squarecross]   (D) [right=of C] {};
		\draw[chain_arrow] let \p1 = (A.two), \p2 = (A.center) in (\x1,\y2) -- (B);
		\draw[chain_arrow] let \p1 = (B.two), \p2 = (B.center) in (\x1,\y2) -- (C);
		\draw[chain_arrow] let \p1 = (C.two), \p2 = (C.center) in (\x1,\y2) -- (D);
	\end{tikzpicture}
\end{center}


\section*{Objetivo}

Podemos crear una \texttt{LinkedList} encadenando nodos simples. 
Pero, en este caso, la lista no es realmente usable.
Tenemos que almacenar el primer elemento en una variable y usar bucles para acceder a cualquier elemento que queramos.
Queremos usar la \texttt{LinkedList} de la misma manera que usamos una \texttt{list} en Python. 
El primer paso es la encapsulación de la funcionalidad de una lista en una clase.
Definiremos métodos en esa clase que proporcionarán una funcionalidad similar a una \texttt{list} normal.

Toma el archivo python proporcionado y amplía la clase para incluir las funcionalidades que faltan.
Puedes echar un vistazo a las diapositivas proporcionadas para obtener algunas pistas sobre cómo implementar cada una de las características.
Puedes intentar programarlas en el orden que prefieras, no es necesario seguir el orden mostrado en este documento.

\textbf{Las funciones deben tener los nombres mostrados en este documento.}

\newpage
\subsection*{Insertar elementos}

\begin{enumerate}
	\item Insertar en la cabeza:
\begin{lstlisting}
lst = LinkedList()
lst.add_head("a")
print(lst)
# Salida: a -> None
lst.add_head("b")
print(lst)
# Salida: b -> a -> None
lst.add_head("c")
print(lst)
# Salida: c -> b -> a -> None
\end{lstlisting}
	\item Insertar al final (cola):
\begin{lstlisting}
lst = LinkedList()
lst.add_tail("a")
print(lst)
# Salida: a -> None
lst.add_tail("b")
print(lst)
# Salida: a -> b -> None
lst.add_tail("c")
print(lst)
# Salida: a -> b -> c -> None
\end{lstlisting}
	\item Insertar en cualquier lugar usando un índice (recuerda que los índices empiezan en 0). El nuevo elemento se insertará después del índice dado:
\begin{lstlisting}
print(lst)
# Salida: a -> b -> c -> None
lst.add_after(1, "d")
print(lst)
# Salida: a -> b -> d -> c -> None
lst.add_after(1, "e")
print(lst)
# Salida: a -> b -> e -> d -> c -> None
\end{lstlisting}
\end{enumerate}

\newpage

\subsection*{Eliminar elementos}

\begin{enumerate}
	\item Eliminar la cabeza:
\begin{lstlisting}
print(lst)
# Salida: a -> b -> c -> None
lst.remove_head()
print(lst)
# Salida: b -> c -> None
lst.remove_head()
print(lst)
# Salida: c -> None
\end{lstlisting}
	\item Eliminar el final (cola):
\begin{lstlisting}
print(lst)
# Salida: a -> b -> c -> None
lst.remove_tail()
print(lst)
# Salida: a -> b -> None
lst.remove_tail()
print(lst)
# Salida: a -> None
\end{lstlisting}
	\item Eliminar en un índice:
\begin{lstlisting}
print(lst)
# Salida: a -> b -> c -> e -> d -> None
lst.remove_index(2)
print(lst)
# Salida: a -> b -> e -> d -> None
lst.remove_index(0)
print(lst)
# Salida: b -> e -> d -> None
\end{lstlisting}
	\item Eliminar un elemento dado:
\begin{lstlisting}
print(lst)
# Salida: a -> b -> c -> e -> d -> None
lst.remove("b")
print(lst)
# Salida: a -> c -> e -> d -> None
lst.remove("e")
print(lst)
# Salida: a -> c -> d -> None
\end{lstlisting}
\end{enumerate}

\newpage

\subsection*{Otras funcionalidades}

\begin{enumerate}
	\item Imprimir la lista de forma amigable en la consola usando \lstinline|print(lst)|
	\item Obtener el tamaño de la lista:
\begin{lstlisting}
print(lst)
# Salida: a -> b -> c -> e -> d -> None
print(lst.size())
# Salida: 5
print(len(lst))
# Salida: 5
\end{lstlisting}
	\item Comprobar si un elemento existe en la lista:
\begin{lstlisting}
print(lst)
# Salida: a -> b -> c -> e -> d -> None
print(lst.contains("a"))
# Salida: True
print(lst.contains("z"))
# Salida: False
print("c" in lst)
# Salida: True
print("x" in lst)
# Salida: False
\end{lstlisting}
	\item Obtener un elemento en un índice:
\begin{lstlisting}
print(lst)
# Salida: a -> b -> c -> e -> d -> None
print(lst.get(1))
# Salida: b
print(lst[2])
# Salida: c
\end{lstlisting}
	\item Modificar un elemento en un índice:
\begin{lstlisting}
print(lst)
# Salida: a -> b -> c -> e -> d -> None
lst.set(1, "z")
print(lst)
# Salida: a -> z -> c -> d -> e -> None
lst[2] = "x"
print(lst)
# Salida: a -> z -> x -> d -> e -> None
\end{lstlisting}
\end{enumerate}

\section*{Evaluando el Rendimiento}
No todas las estructuras de lista se comportan igual. Algunas pueden ser más eficientes en una operación, otras en otra y algunas pueden intentar encontrar un equilibrio. Vamos a comparar el rendimiento de la estructura \lstinline|list| en Python, con nuestra propia \lstinline|LinkedList| y una implementación de lista enlazada proporcionada por el módulo \lstinline|collections|, el \lstinline|deque|. Para ver qué operación tiene mejor rendimiento, simplemente podemos usar el tiempo que tarda un ordenador en ejecutarla.

\begin{attachedlisting}[time]
import time

inicio = time.perf_counter_ns()
x = (12 ** 124) ** 0.5
fin = time.perf_counter_ns()

print(f"Tiempo: {(fin-inicio)/1e9:.10f} segundos")
\end{attachedlisting}

Echa un vistazo al código, ¿cuánto tiempo necesita para realizar el cálculo? Ejecuta el código varias veces, ¿el tiempo es siempre el mismo o varía? ¿Qué sucede si eliminamos el código de la línea 4? ¿Necesita algo de tiempo para ejecutar simplemente las sentencias para obtener el tiempo en nanosegundos? ¿Afectará a nuestras mediciones?

Podemos resolver la mayoría de estos problemas simplemente ejecutando el mismo fragmento de código varias veces. Al ejecutar el mismo fragmento de código, digamos 1000 veces, reducimos el efecto de las variaciones y el tiempo perdido obteniendo los contadores.

\begin{attachedlisting}[time]
import time

inicio = time.perf_counter_ns()
for _ in range(1000):
    x = (12 ** 124) ** 0.5
fin = time.perf_counter_ns()

print(f"Tiempo: {(fin-inicio)/1e9:.10f} segundos")
\end{attachedlisting}

Ejecuta esto varias veces, ¿el tiempo sigue variando? Quizás en lugar de ejecutarlo 1000 veces, aumenta el valor a 1000000. ¿Qué sucede? Para obtener datos más consistentes necesitamos ejecutar el código bastantes veces.

Utiliza la misma estructura para obtener el tiempo que una \lstinline|list|, \lstinline|LinkedList| y un \lstinline|deque| necesitan para las siguientes operaciones, y rellena la tabla con el tiempo que requiere para la cantidad de operaciones. Para el primer conjunto de operaciones empezaremos con una lista vacía y añadiremos $n$ elementos en un bucle. Para las operaciones de obtención (get) necesitamos algunos datos en las listas, así que podemos simplemente añadir algunos puntos de datos ficticios a las listas y luego intentar obtenerlos.

\newpage
Número de iteraciones: \rule{3cm}{0.15mm}\\

\begin{table}[h!bt]
\renewcommand{\arraystretch}{1.4}
\centering
\begin{tabular}{m{0.3\textwidth}m{0.2\textwidth}m{0.2\textwidth}m{0.2\textwidth}}
\toprule
Operación & \lstinline|list| & \lstinline|LinkedList| & \lstinline|deque| \\
\midrule
Insertar en la cabeza & & &\\
Insertar al final & & &\\
Insertar en índice & & &\\
\midrule
Obtener elemento cabeza & & &\\
Obtener elemento final & & &\\
Obtener en índice & & &\\
\bottomrule
\end{tabular}
\end{table}

\FloatBarrier

\newpage

\section*{Otras mejoras a una lista}

\begin{enumerate}
	\item Hemos creado una lista donde podemos insertar un elemento en cualquier posición que queramos. Ahora queremos tener una lista que esté ordenada. Crea una clase \texttt{SortedLinkedList} donde solo tengamos un método que permita la inserción de elementos \texttt{SortedLinkedList.add(data)}. El método insertará el elemento en la posición correcta para que todos los elementos de la lista antes del elemento sean menores, y todos los elementos después sean mayores. Usa los operandos $<, >, <=, >=$ para comparar los elementos de la lista con el elemento a insertar y encontrar la posición donde tiene que ser añadido.
	\item Hemos creado una lista donde cada nodo solo referencia al siguiente. Esto crea una lista donde añadir o eliminar la cabeza de la lista es realmente simple y muy eficiente. Podemos ampliar esta funcionalidad creando una \texttt{DoublyLinkedList}; en este caso, cada nodo no solo apuntará al siguiente elemento de la lista, sino que también contendrá la referencia al nodo anterior de la lista. En este caso la clase no solo mantendrá una referencia a la cabeza, sino también al final (tail) de la lista. Creando efectivamente una estructura con dos listas yendo en direcciones opuestas. Usar esta operación en el final es tan eficiente como la operación en la cabeza.
    \begin{figure}[h!bt]
	\begin{center}
	\shorthandoff{>}
	\begin{tikzpicture}[
		list/.style={
						very thick, rectangle split,
						rectangle split parts=3, draw,
						rectangle split horizontal, minimum size=18pt,
						inner sep=5pt, text=black,
						rectangle split part fill={blue!20, red!20, blue!20}
					},
		->, start chain, very thick
		]
		
		\node[list,on chain] (A) {\nodepart{second} cabeza};
		\node[list,on chain] (B) {\nodepart{second} medio};
		\node[list,on chain] (C) {\nodepart{second} cola};
		
		\node[squarecross]   (D) [right=of C] {};
		\node[squarecross]   (E) [left= of A] {};
		
		\path[*->] let \p1 = (A.three), \p2 = (A.center) in (\x1,\y2) edge [bend left] ($(B.west)+(0,0.2)$);
		\path[*->] let \p1 = (B.three), \p2 = (B.center) in (\x1,\y2) edge [bend left] ($(C.west)+(0,0.2)$);
		\draw[*->] let \p1 = (C.three), \p2 = (C.center) in (\x1,\y2) -- (D);
		
		\draw[*->] ($(A.one)+(0.2,0.1)$) -- (E);
		\path[*->] ($(B.one)+(0.1,0.1)$) edge [bend left] ($(A.east)+(0,-0.05)$);
		\path[*->] ($(C.one)+(0.1,0.1)$) edge [bend left] ($(B.east)+(0,-0.05)$);
	\end{tikzpicture}
	\end{center}
    \end{figure}
\end{enumerate}

\end{document}